\section{September 21, 2026}

\subsection{Spanning collections and dimension}

\begin{proposition}\label{prop:spanning-dim-is-basis}
  Suppose that $V$ is finite-dimensional. Every collection of
  $\dim V$ vectors that spans $V$ is a basis of $V$.
\end{proposition}

\begin{proof}
  Write $\dim V=n$, and suppose that $v_1,\ldots,v_n$ spans $V$.
  This collection can be reduced to a basis by removing redundant
  vectors. However, every basis of $V$ has exactly $n$ vectors, so no
  vectors can be removed. Thus $v_1,\ldots,v_n$ is already a basis of
  $V$.
\end{proof}

\subsection{Dimension of a sum of subspaces}

\begin{proposition}\label{prop:dimension-of-sum}
  If $V_1$ and $V_2$ are subspaces of a finite-dimensional vector space,
  then
  \[
    \dim(V_1+V_2)=\dim V_1+\dim V_2-\dim(V_1\cap V_2).
  \]
\end{proposition}

\begin{proof}
  Let $v_1,\ldots,v_m$ be a basis of $V_1\cap V_2$, so that
  $\dim(V_1\cap V_2)=m$. This collection is linearly independent in
  $V_1$, so the basis extension theorem gives a basis
  \[
    v_1,\ldots,v_m,u_1,\ldots,u_j
  \]
  of $V_1$. Similarly, extend $v_1,\ldots,v_m$ to a basis
  \[
    v_1,\ldots,v_m,w_1,\ldots,w_k
  \]
  of $V_2$. Hence
  \[
    \dim V_1=m+j
    \qquad\text{and}\qquad
    \dim V_2=m+k.
  \]
  We will show that
  \[
    v_1,\ldots,v_m,u_1,\ldots,u_j,w_1,\ldots,w_k
  \]
  is a basis of $V_1+V_2$.

  All the displayed vectors belong to $V_1+V_2$, so their span is
  contained in $V_1+V_2$. Conversely, every vector in $V_1+V_2$ is a
  sum of a vector in $V_1$ and a vector in $V_2$. Expressing each
  summand in the corresponding basis shows that the combined
  collection spans $V_1+V_2$.

  To prove linear independence, suppose that
  \begin{equation}\label{eq:dimension-sum-dependence}
    a_1v_1+\cdots+a_mv_m
    +b_1u_1+\cdots+b_ju_j
    +c_1w_1+\cdots+c_kw_k=0.
  \end{equation}
  Rearranging gives
  \[
    c_1w_1+\cdots+c_kw_k
    =-a_1v_1-\cdots-a_mv_m-b_1u_1-\cdots-b_ju_j.
  \]
  The right-hand side belongs to $V_1$, while the left-hand side
  belongs to $V_2$. Therefore
  \[
    c_1w_1+\cdots+c_kw_k\in V_1\cap V_2.
  \]
  Since $v_1,\ldots,v_m$ is a basis of the intersection, there are
  scalars $d_1,\ldots,d_m$ such that
  \[
    c_1w_1+\cdots+c_kw_k=d_1v_1+\cdots+d_mv_m.
  \]
  Thus
  \[
    c_1w_1+\cdots+c_kw_k-d_1v_1-\cdots-d_mv_m=0.
  \]
  The collection $v_1,\ldots,v_m,w_1,\ldots,w_k$ is a basis of $V_2$
  and hence linearly independent. Consequently, all the $c$'s and
  $d$'s are zero. Equation~\eqref{eq:dimension-sum-dependence} now
  becomes
  \[
    a_1v_1+\cdots+a_mv_m+b_1u_1+\cdots+b_ju_j=0.
  \]
  Because $v_1,\ldots,v_m,u_1,\ldots,u_j$ is a basis of $V_1$, all
  the $a$'s and $b$'s are also zero. This proves linear independence.

  The combined collection is therefore a basis of $V_1+V_2$, and
  \begin{align*}
    \dim(V_1+V_2)
      &=m+j+k \\
      &=(m+j)+(m+k)-m \\
      &=\dim V_1+\dim V_2-\dim(V_1\cap V_2).
  \end{align*}
\end{proof}

\subsection{Linear maps on finite-dimensional spaces}

\begin{theorem}[Fundamental theorem of linear maps]
  \label{thm:fundamental-linear-maps}
  Suppose that $V$ is finite-dimensional and that
  $T\in\mathcal L(V,W)$. Then $\operatorname{im}T$ is
  finite-dimensional and
  \[
    \dim V=\dim\ker T+\dim\operatorname{im}T.
  \]
\end{theorem}

Recall that
\[
  \ker T=\{v\in V:T(v)=0\}
  \qquad\text{and}\qquad
  \operatorname{im}T=\{T(v):v\in V\}.
\]
Before proving the theorem, we connect the kernel and image with
injectivity and surjectivity, beginning with injectivity.

\begin{proposition}\label{prop:injective-kernel}
  Let $T\in\mathcal L(V,W)$. Then $T$ is injective if and only if
  $\ker T=\{0\}$.
\end{proposition}

\begin{proof}
  First note that every linear map sends zero to zero, since
  \[
    T(0)=T(0+0)=T(0)+T(0)
  \]
  implies $T(0)=0$. This does not require injectivity. In particular,
  $\{0\}\subseteq\ker T$.

  Suppose that $T$ is injective. If $v\in\ker T$, then
  \[
    T(v)=0=T(0).
  \]
  Injectivity gives $v=0$, so $\ker T\subseteq\{0\}$. Therefore
  $\ker T=\{0\}$.

  Conversely, suppose that $\ker T=\{0\}$. Let $v_1,v_2\in V$ satisfy
  $T(v_1)=T(v_2)$. By linearity,
  \[
    T(v_1-v_2)=T(v_1)-T(v_2)=0.
  \]
  Hence $v_1-v_2\in\ker T=\{0\}$, so $v_1=v_2$. Thus $T$ is
  injective.
\end{proof}
